Now consider the reductive \(k\)-group
\[
  H=G\times_k\mathbb G_{m,k}.
\]
There is an \(H\)-stable decomposition
\[
  \Lie(H/k)=\Lie(G/k)\oplus k.
\]
For every \(H\)-module \(V\),
\[
  H^i(H,V)=H^i\bigl(G,H^0(\mathbb G_{m,k},V)\bigr).
\]
Indeed, the \(H^i(H,-)\) are the derived functors of
\[
  H^0(H,-)
  =H^0\bigl(G,H^0(\mathbb G_{m,k},-)\bigr),
\]
and the functor \(H^0(\mathbb G_{m,k},-)\) is exact
(cf.\ \SGARefLink{sga3:i:theorem:5-3-1}{Exp.~I, 5.3.1} and \SGARefLink{sga3:i:theorem:5-3-3}{5.3.3}). In particular, for every \(i\),
\[
  H^i(H,\Lie(H/k))
  =
  H^i(G,\Lie(G/k))\oplus H^i(G,k).
\]
Applying \SGARefLink{sga3:xxiv:1.13:item:ii}{1.13(ii)} to the reductive group \(H\) gives
\(H^2(G,k)=0\).

\SGARefTarget{sga3:xxiv:1.15.1}
\medskip\noindent\textbf{Corollary 1.15.1.} ---
\footnote{Editor's note: This corollary, drawn from remarks of
O.~Gabber, has been added to make \SGARefLink{sga3:xxiv:1.13:item:i}{1.13(i)} more precise.}
{\itshape Let \(k\) be a field, \(G\) a reductive \(k\)-group, \(Z\)
its center, and \(\mathscr R=(M,M^*,R,R^*)\) the type of \(G\). Then
\[
  H^1(G,\Lie(G/k))
  \simeq
  \operatorname{Ext}^1_{\mathbb Z}(M/\Gamma_0(R),k).
\]
In particular, \(H^1(G,\Lie(G/k))=0\) if and only if \(Z\) is smooth
over \(k\).}

Let \(\mathfrak g\), \(\mathfrak z\), and
\(\mathfrak g_{\mathrm{ad}}\) be the Lie algebras of \(G\), \(Z\), and
\(G_{\mathrm{ad}}=G/Z\), respectively. It follows from \SGARefLink{sga3:xxiv:1.15}{1.15} and its
proof that
\[
  H^1(G,\mathfrak g)
  =
  H^1(G_{\mathrm{ad}},\mathfrak g)
  \simeq
  H^1(G_{\mathrm{ad}},\mathfrak g/\mathfrak z).
\]
Set \(C=\operatorname{Coker}(\mathfrak g\to
\mathfrak g_{\mathrm{ad}})\). There is an exact sequence
\[
  0\longrightarrow\mathfrak g/\mathfrak z
  \longrightarrow\mathfrak g_{\mathrm{ad}}
  \longrightarrow C\longrightarrow0.
\]
Since \(H^1(G_{\mathrm{ad}},\mathfrak g_{\mathrm{ad}})=0\) by \SGARefLink{sga3:xxiv:1.13}{1.13}
and
\[
  H^0(G_{\mathrm{ad}},\mathfrak g_{\mathrm{ad}})=0
\]
(cf.\ \SGARefLink{sga3:ii:theorem:5-2-3}{Exp.~II, 5.2.3}), we obtain
\[
  H^1(G_{\mathrm{ad}},\mathfrak g/\mathfrak z)
  \simeq H^0(G_{\mathrm{ad}},C).
\]
To calculate the right-hand side, we may assume \(k\) algebraically
closed. Let \((G,T,M,R)\) be a splitting of \(G\). Then
\[
  Z=D_k(M/\Gamma_0(R))
\]
and \(C\) is identified with
\[
  \operatorname{Coker}\!\left(
    \operatorname{Hom}_{\mathbb Z}(M,k)
    \longrightarrow
    \operatorname{Hom}_{\mathbb Z}(\Gamma_0(R),k)
  \right)
  \simeq
  \operatorname{Ext}^1_{\mathbb Z}(M/\Gamma_0(R),k),
\]
with trivial \(G_{\mathrm{ad}}\)-action. The corollary follows, since
\(Z\) is not smooth over \(k\) exactly when
\(\operatorname{char}(k)=p>0\) and \(M/\Gamma_0(R)\) has
\(p\)-torsion (\SGARefLink{sga3:ix:proposition:2:1}{Exp.~IX, 2.1}).

\SGARefTarget{sga3:xxiv:1.16}
\medskip\noindent\textbf{Definition 1.16.} ---
{\itshape Let \(S\) be a scheme and \(G\) a reductive \(S\)-group. A
\emph{form of \(G\) over \(S\)} is an \(S\)-group scheme \(G'\)
locally isomorphic to \(G\) for the \textup{(fpqc)} topology. It is
equivalent to require that \(G'\) be locally isomorphic to \(G\) for
the étale topology (cf.\ \SGARefLink{sga3:XXIII:5.6}{Exp.~XXIII, 5.6}), or that \(G'\) be a
reductive \(S\)-group of the same type as \(G\) at every point of
\(S\).}

\SGARefTarget{sga3:xxiv:1.17}
\medskip\noindent\textbf{Corollary 1.17.} ---
{\itshape Let \(S\) be a scheme and \(G\) a reductive \(S\)-group.
\begin{enumerate}
\item[(i)] The functor
\[
  G'\longmapsto
  \underline{\operatorname{Isom}}_{S\text{-}\mathrm{gr}}(G,G')
\]
is an equivalence between the category of forms of \(G\) over \(S\)
and the category of principal homogeneous bundles under
\(\underline{\Aut}_{S\text{-}\mathrm{gr}}(G)\).
\item[(ii)] If \(S'\to S\) is a covering morphism, forms of \(G\)
trivialized by \(S'\) correspond to bundles trivialized by \(S'\).
\item[\SGAInlineRefTarget{sga3:xxiv:1.17:item:iii}(iii)] Every principal homogeneous sheaf under
\(\underline{\Aut}_{S\text{-}\mathrm{gr}}(G)\) is representable and
quasi-isotrivial, that is, locally trivial for the étale topology.
\end{enumerate}}

The first assertion is formal in the category of sheaves, for example
for the \textup{(fpqc)} topology. Every sheaf locally isomorphic to
\(G\) for the \textup{(fpqc)} topology is representable, since \(G\)
is affine over \(S\), and is locally isomorphic to \(G\) for the
étale topology. Finally, for every form \(G'\) of \(G\), the
\(S\)-sheaf
\(\underline{\operatorname{Isom}}_{S\text{-}\mathrm{gr}}(G,G')\) is
representable by \SGARefLink{sga3:xxiv:1.8}{1.8}. The corollary follows.

\SGARefTarget{sga3:xxiv:1.18}
\medskip\noindent\textbf{Corollary 1.18.} ---
{\itshape The set of isomorphism classes of forms of the reductive
group \(G\) over \(S\) is isomorphic to
\[
  H^1\!\left(S,\underline{\Aut}_{S\text{-}\mathrm{gr}}(G)\right)
  =
  H^1_{\mathrm{\acute et}}\!
    \left(S,\underline{\Aut}_{S\text{-}\mathrm{gr}}(G)\right)
  =
  \Fib\!\left(S,\underline{\Aut}_{S\text{-}\mathrm{gr}}(G)\right).
\]
If \(S'\to S\) is a covering morphism, the subset consisting of forms
trivialized by \(S'\) is isomorphic to
\[
  H^1\!\left(S'/S,\underline{\Aut}_{S\text{-}\mathrm{gr}}(G)\right).
\]}

\SGARefTarget{sga3:xxiv:1.19}
\medskip\noindent\textbf{Corollary 1.19.} ---
{\itshape Let \(S\) be a scheme and \(\mathscr R\) a pinned reduced root
datum such that \(\operatorname{\acute{E}p}_S(\mathscr R)\)
exists,\footnote{Editor's note: Cf.\ XXIII, \SGARefLink{sga3:XXIII:5.11}{Definition 5.11}.} a
condition which is automatically satisfied by Exp.~XXV. Put
\[
  A_S(\mathscr R)
  =
  \underline{\Aut}_{S\text{-}\mathrm{gr}}
    (\operatorname{\acute{E}p}_S(\mathscr R))
  =
  \operatorname{ad}(\operatorname{\acute{E}p}_S(\mathscr R))
  \cdot E(\mathscr R)_S.
\]
\begin{enumerate}
\item[\SGAInlineRefTarget{sga3:xxiv:1.19:item:i}(i)] The set of isomorphism classes of reductive \(S\)-groups
of type \(\mathscr R\) (\SGARefLink{sga3:xxii:def:2.7}{Exp.~XXII, 2.7}) is isomorphic, by
\SGARefLink{sga3:XXIII:5.12}{Exp.~XXIII, 5.12}, to
\[
  H^1(S,A_S(\mathscr R))
  =
  H^1_{\mathrm{\acute et}}(S,A_S(\mathscr R))
  =
  \Fib(S,A_S(\mathscr R)).
\]
\item[\SGAInlineRefTarget{sga3:xxiv:1.19:item:ii}(ii)] If \(S'\to S\) is a covering morphism, the subset
consisting of classes of groups splittable over \(S'\) is isomorphic
to \(H^1(S'/S,A_S(\mathscr R))\).
\end{enumerate}}

\SGARefTarget{sga3:xxiv:1.20}
\medskip\noindent\textbf{Remark 1.20.} ---
With the preceding notation, every reductive \(S\)-group \(G\) of
type \(\mathscr R\) has a canonically associated right principal
homogeneous bundle under \(A_S(\mathscr R)\):
\[
  \underline{\operatorname{Isom}}_{S\text{-}\mathrm{gr}}
    (\operatorname{\acute{E}p}_S(\mathscr R),G)=P.
\]
The scheme \(P\) may be interpreted as the ``scheme of pinnings of
\(G\) of type \(\mathscr R\)'' (cf.\ \SGARefLink{sga3:xxiv:1.0}{1.0}). It is also a principal
homogeneous bundle on the left under
\(\underline{\Aut}_{S\text{-}\mathrm{gr}}(G)\), a structure visible
immediately from the description above.

\SGARefTarget{sga3:xxiv:1.21}
\medskip\noindent\textbf{Proposition 1.21.} ---
{\itshape Let \(S\) be a henselian local scheme and \(s\) its closed
point. The functor
\[
  G\longmapsto G_s
\]
induces a bijection from the set of isomorphism classes of reductive
\(S\)-groups onto the set of isomorphism classes of reductive
\(\kappa(s)\)-groups. In particular, for every reductive \(S\)-group
\(G\), there exists a finite étale surjective morphism \(S'\to S\)
such that \(G_{S'}\) is splittable.}

Using the existence of \(\operatorname{\acute{E}p}_S(\mathscr R)\)
(Exp.~XXV), it is enough to prove that, with
\(H=A_S(\mathscr R)\), the canonical map
\[
  \Fib(S,H)\longrightarrow\Fib(\kappa(s),H_s)
\]
is bijective, and that every element of \(\Fib(S,H)\) has the property
just stated. Every finite subset of \(H\) is contained in an affine
open subscheme: this is clear for a constant group, and \(H\) is
affine over a constant group. The result proved in the Appendix,
\SGARefLink{sga3:xxiv:proposition:8-1}{8.1}, therefore applies.

\SGARefTarget{sga3:xxiv:section:2}
\section*{2. Automorphisms and Subgroups}

Introduce the following notation. If
\[
  \mathcal H=\underline{\Aut}_{S\text{-}\mathrm{gr}}(G)
\]
and \(X\) is a subfunctor of \(G\), set
\[
  \underline{\Aut}_{S\text{-}\mathrm{gr}}(G,X)
  =\operatorname{Norm}_{\mathcal H}(X),
  \qquad
  \underline{\Aut}_{S\text{-}\mathrm{gr}}(G,\operatorname{id}_X)
  =\underline{\Centr}_{\mathcal H}(X).
\]
If \(Y\) is a second subfunctor of \(G\), define
\[
  \underline{\Aut}_{S\text{-}\mathrm{gr}}(G,X,Y)
  =
  \underline{\Aut}_{S\text{-}\mathrm{gr}}(G,X)
  \cap
  \underline{\Aut}_{S\text{-}\mathrm{gr}}(G,Y).
\]
If \(G'\) is a second \(S\)-group and \(X'\) a subfunctor of \(G'\),\footnote{Translator's note: The French text prints \(G\) here; \(G'\) is required by the typing of the definition.}
write
\[
  \underline{\operatorname{Isom}}_{S\text{-}\mathrm{gr}}
    (G,X;G',X')
\]
for the subfunctor of
\(\underline{\operatorname{Isom}}_{S\text{-}\mathrm{gr}}(G,G')\)
defined, for every \(S'\to S\), by
\[
  \left\{
    u\in
    \underline{\operatorname{Isom}}_{S\text{-}\mathrm{gr}}(G,G')(S')
    \,\middle|\,
    u(X_{S'})=X'_{S'}
  \right\}.
\]
The notation
\(\underline{\operatorname{Isom}}_{S\text{-}\mathrm{gr}}
(G,X,Y;G',X',Y')\), and so on, is defined similarly.\footnote{Editor's
note: The preceding definitions have been made explicit. The original
only said that
\(\underline{\Aut}_{S\text{-}\mathrm{gr}}(G,X,Y),\ldots\) and
\(\underline{\operatorname{Isom}}_{S\text{-}\mathrm{gr}}
(G,X;G',X'),\ldots\) were defined similarly.}

\SGARefTarget{sga3:xxiv:2.1}
\medskip\noindent\textbf{Proposition 2.1.} ---
{\itshape Let \(S\) be a scheme, \(G\) a reductive \(S\)-group, and
\(T\) a maximal torus of \(G\) (respectively, \(B\) a Borel subgroup
of \(G\), or \(B\supset T\) a Killing couple of \(G\)). Let
\(T_{\mathrm{ad}}\) (respectively, \(B_{\mathrm{ad}}\)) denote the
corresponding maximal torus (respectively, Borel subgroup) of
\(\operatorname{ad}(G)\):
\[
  B_{\mathrm{ad}}
  \simeq B/\underline{\Centr}(G)
  =B/\underline{\Centr}(B),
  \qquad
  T_{\mathrm{ad}}\simeq T/\underline{\Centr}(G).
\]
Then
\(\underline{\Aut}_{S\text{-}\mathrm{gr}}(G,T)\)
(respectively,
\(\underline{\Aut}_{S\text{-}\mathrm{gr}}(G,B)\), or
\(\underline{\Aut}_{S\text{-}\mathrm{gr}}(G,B,T)\))
is represented by a closed subscheme of
\(\underline{\Aut}_{S\text{-}\mathrm{gr}}(G)\), smooth over \(S\),
and the exact sequence of \SGARefLink{sga3:xxiv:1.3}{Theorem 1.3} induces the exact sequences
\[
  1\longrightarrow
  \operatorname{Norm}_{\operatorname{ad}(G)}(T_{\mathrm{ad}})
  \longrightarrow
  \underline{\Aut}_{S\text{-}\mathrm{gr}}(G,T)
  \longrightarrow
  \underline{\Autext}(G)
  \longrightarrow1,
\]
\[
  1\longrightarrow B_{\mathrm{ad}}
  \longrightarrow
  \underline{\Aut}_{S\text{-}\mathrm{gr}}(G,B)
  \longrightarrow
  \underline{\Autext}(G)
  \longrightarrow1,
\]
\[
  1\longrightarrow T_{\mathrm{ad}}
  \longrightarrow
  \underline{\Aut}_{S\text{-}\mathrm{gr}}(G,B,T)
  \longrightarrow
  \underline{\Autext}(G)
  \longrightarrow1.
\]}

By descent of closed subschemes,\footnote{Editor's note: Cf.\ SGA~1,
VIII, 1.9.} we reduce immediately to the case in which \(G\) is
pinned and \(B\supset T\) is its canonical Killing couple
(cf.\ \SGARefLink{sga3:xxii:cor:5.5.5:item:iv}{Exp.~XXII, 5.5.5(iv)}). Since the group \(E\) of \SGARefLink{sga3:xxiv:1.2}{1.2} normalizes
\(B\) and \(T\), the result follows at once from the normalization
theorems in \(\operatorname{ad}(G)\) (\SGARefLink{sga3:xxii:cor:5.3.12}{Exp.~XXII, 5.3.12} and 5.6.1).


Using now the conjugacy theorems (cf.\ \SGARefLink{sga3:XXIII:5.12}{Exp.~XXIII, 5.12}), and arguing
as in \S1, one obtains:

\SGARefTarget{sga3:xxiv:2.2}
\medskip\noindent\textbf{Corollary 2.2.} ---
{\itshape Let \(S\) be a scheme, let \(G\) and \(G'\) be two reductive
\(S\)-groups having the same type at every point, and let
\(B\supset T\) (respectively, \(B'\supset T'\)) be a Killing couple
of \(G\) (respectively, \(G'\)).

\textup{\SGARefLink{sga3:xxiv:1.19:item:i}{(i)}} The \(S\)-functor
\[
  \underline{\operatorname{Isom}}_{S\text{-}\mathrm{gr}}
    (G,T;G',T')
\]
is represented by a closed smooth subscheme of
\(\underline{\operatorname{Isom}}_{S\text{-}\mathrm{gr}}(G,G')\)
which is a principal homogeneous space under
\(\underline{\Aut}_{S\text{-}\mathrm{gr}}(G,T)\).
Moreover,
\(\operatorname{Norm}_{\operatorname{ad}(G)}(T_{\mathrm{ad}})\)
acts freely on this scheme, and there is a canonical isomorphism
\[
  \underline{\operatorname{Isom}}_{S\text{-}\mathrm{gr}}
    (G,T;G',T')
  \big/
  \operatorname{Norm}_{\operatorname{ad}(G)}(T_{\mathrm{ad}})
  \simeq
  \underline{\operatorname{Isomext}}(G,G').
\]

\textup{\SGARefLink{sga3:xxiv:1.19:item:ii}{(ii)}} The \(S\)-functor
\[
  \underline{\operatorname{Isom}}_{S\text{-}\mathrm{gr}}
    (G,B;G',B')
\]
is represented by a closed smooth subscheme of
\(\underline{\operatorname{Isom}}_{S\text{-}\mathrm{gr}}(G,G')\)
which is a principal homogeneous space under
\(\underline{\Aut}_{S\text{-}\mathrm{gr}}(G,B)\).
Moreover, \(B_{\mathrm{ad}}\) acts freely on this scheme, and there
is a canonical isomorphism
\[
  \underline{\operatorname{Isom}}_{S\text{-}\mathrm{gr}}
    (G,B;G',B')
  \big/B_{\mathrm{ad}}
  \simeq
  \underline{\operatorname{Isomext}}(G,G').
\]

\textup{\SGARefLink{sga3:xxiv:1.17:item:iii}{(iii)}} The \(S\)-functor
\[
  \underline{\operatorname{Isom}}_{S\text{-}\mathrm{gr}}
    (G,B,T;G',B',T')
\]
is represented by a closed smooth subscheme of
\(\underline{\operatorname{Isom}}_{S\text{-}\mathrm{gr}}(G,G')\),
which is a principal homogeneous space under
\(\underline{\Aut}_{S\text{-}\mathrm{gr}}(G,B,T)\).
Moreover, \(T_{\mathrm{ad}}\) acts freely on this scheme, and there
is a canonical isomorphism
\[
  \underline{\operatorname{Isom}}_{S\text{-}\mathrm{gr}}
    (G,B,T;G',B',T')
  \big/T_{\mathrm{ad}}
  \simeq
  \underline{\operatorname{Isomext}}(G,G').
\]}

Arguing once again as in \S1, one obtains:

\SGARefTarget{sga3:xxiv:2.3}
\medskip\noindent\textbf{Corollary 2.3.} ---
{\itshape Let \(S\) be a scheme, let \(G\) be a reductive \(S\)-group,
and let \(B\supset T\) be a Killing couple of \(G\). The functor}
\[
  (G',T')
  \longmapsto
  \underline{\operatorname{Isom}}_{S\text{-}\mathrm{gr}}
    (G,T;G',T'),
\]
{\itshape respectively}
\[
  (G',B')
  \longmapsto
  \underline{\operatorname{Isom}}_{S\text{-}\mathrm{gr}}
    (G,B;G',B'),
\]
{\itshape respectively}
\[
  (G',B',T')
  \longmapsto
  \underline{\operatorname{Isom}}_{S\text{-}\mathrm{gr}}
    (G,B,T;G',B',T'),
\]
{\itshape is an equivalence between the category of pairs
\((G',T')\) (respectively, pairs \((G',B')\), respectively, triples
\((G',B',T')\)), where \(G'\) is a form of \(G\) and \(T'\) is a
maximal torus of \(G'\) (respectively, \(B'\) is a Borel subgroup of
\(G'\), respectively, \(B'\supset T'\) is a Killing couple of
\(G'\)), and the category of principal homogeneous bundles under
the \(S\)-group \(H\), where
\(H=\underline{\Aut}_{S\text{-}\mathrm{gr}}(G,T)\)
(respectively,
\(H=\underline{\Aut}_{S\text{-}\mathrm{gr}}(G,B)\),
respectively,
\(H=\underline{\Aut}_{S\text{-}\mathrm{gr}}(G,B,T)\)).
Moreover, every principal homogeneous sheaf under \(H\) is
representable and quasi-isotrivial, so that}
\[
  H^1(S,H)=H^1_{\mathrm{\acute et}}(S,H)=\Fib(S,H).
\]


\SGARefTarget{sga3:xxiv:2.4}
\medskip\noindent\textbf{Remark 2.4.} ---
Under the hypotheses of \SGARefLink{sga3:xxiv:2.2}{2.2}, the morphism denoted set-theoretically
by \(u\mapsto u(T)\) (respectively, \(u\mapsto u(B)\), respectively,
\(u\mapsto(u(B),u(T))\)) induces an isomorphism
\[
  \underline{\operatorname{Isom}}_{S\text{-}\mathrm{gr}}(G,G')
  \big/
  \underline{\Aut}_{S\text{-}\mathrm{gr}}(G,T)
  \simeq \operatorname{Tor}(G'),
\]
respectively,
\[
  \underline{\operatorname{Isom}}_{S\text{-}\mathrm{gr}}(G,G')
  \big/
  \underline{\Aut}_{S\text{-}\mathrm{gr}}(G,B)
  \simeq \operatorname{Bor}(G'),
\]
respectively,
\[
  \underline{\operatorname{Isom}}_{S\text{-}\mathrm{gr}}(G,G')
  \big/
  \underline{\Aut}_{S\text{-}\mathrm{gr}}(G,B,T)
  \simeq \operatorname{Kil}(G').
\]
The proof is immediate: it suffices to prove the assertion locally
for \textup{(fpqc)}, so one may assume \(G\simeq G'\), and the result
then reduces to \SGARefLink{sga3:xxii:cor:5.8.3:item:iii}{Exp.~XXII, 5.8.3(iii)}.

\SGARefTarget{sga3:xxiv:2.5}
\medskip\noindent\textbf{Remark 2.5.} ---
The preceding results are immediately interpreted in terms of
reduction of the structure group. If \(G'\) is a form of \(G\),
corresponding to the principal bundle
\(\underline{\operatorname{Isom}}_{S\text{-}\mathrm{gr}}(G,G')\),
to give a reduction of the structure group of this bundle to
\(\underline{\Aut}_{S\text{-}\mathrm{gr}}(G,T)\) amounts to giving
a maximal torus \(T'\) of \(G'\). The indicated bijections are,
on the one hand, that of \SGARefLink{sga3:xxiv:2.4}{2.4} and, on the other hand, the map
\[
  T'\longmapsto
  \underline{\operatorname{Isom}}_{S\text{-}\mathrm{gr}}
    (G,T;G',T').
\]
The same applies to Borel subgroups and Killing couples.

\SGARefTarget{sga3:xxiv:2.6}
\medskip\noindent\textbf{Proposition 2.6.} ---
{\itshape Let \(S\) be a scheme, let \(G\) and \(G'\) be two reductive
\(S\)-groups having the same type at every point, and let \(T\)
(respectively, \(T'\)) be a maximal torus of \(G\) (respectively,
\(G'\)). Then \(T_{\mathrm{ad}}\) acts freely on
\(\underline{\operatorname{Isom}}_{S\text{-}\mathrm{gr}}
(G,T;G',T')\), and the quotient}
\[
  P=
  \underline{\operatorname{Isom}}_{S\text{-}\mathrm{gr}}
    (G,T;G',T')
  \big/T_{\mathrm{ad}}
\]
{\itshape is representable; it is a principal homogeneous bundle under}
\[
  A=
  \underline{\Aut}_{S\text{-}\mathrm{gr}}(G,T)
  \big/T_{\mathrm{ad}},
\]
{\itshape where \(A\) is represented by a twisted constant \(S\)-scheme
and is an extension of \(\underline{\Autext}(G)\) by}
\[
  W_{\operatorname{ad}(G)}(T_{\mathrm{ad}})
  =
  \operatorname{Norm}_{\operatorname{ad}(G)}(T_{\mathrm{ad}})
  \big/T_{\mathrm{ad}}.
\]
{\itshape Moreover, if \(A\) is made to act on \(T\) in the evident
way, the bundle associated with \(P\) is precisely \(T'\).}\footnote{Editor's
note: Consequently, if \(P\) has a section, then \(T\) and \(T'\)
are isomorphic; in particular, \(T_P\) and \(T'_P\) are isomorphic.
See the addition \SGARefLink{sga3:xxiv:proposition:4-5-1}{4.5.1}, where this observation is developed in the
case \(G=\operatorname{\acute{E}p}_S(\mathscr R)\).}

The first part of the proposition follows immediately from the
preceding results. To prove the second, observe that there is an
evident morphism
\[
  P\times_S T\longrightarrow T',
  \qquad (u,t)\longmapsto u(t).
\]
To show that, after passage to the quotient by \(A\), it induces an
isomorphism, one may once again assume
\((G,T)\simeq(G',T')\), in which case the assertion is immediate.

In an entirely analogous way, one has:

\SGARefTarget{sga3:xxiv:2.7}
\medskip\noindent\textbf{Proposition 2.7.} ---
{\itshape Let \(S\) be a scheme, let \(G\) and \(G'\) be two reductive
\(S\)-groups having the same type at every point, and let
\(B\supset T\) (respectively, \(B'\supset T'\)) be a Killing couple
of \(G\) (respectively, \(G'\)). If
\[
  \underline{\Aut}_{S\text{-}\mathrm{gr}}(G,B,T)
  \big/T_{\mathrm{ad}}
  \simeq \underline{\Autext}(G)
\]
is made to act on \(T\) in the evident way, the bundle associated
with \(\underline{\operatorname{Isomext}}(G,G')\) is precisely
\(T'\).}

\SGARefTarget{sga3:xxiv:2.8}
\medskip\noindent\textbf{Corollary 2.8.} ---
{\itshape Let \(G\) and \(G'\) be two reductive \(S\)-groups which are
inner twisted forms of one another, and let \(B\supset T\)
(respectively, \(B'\supset T'\)) be a Killing couple of \(G\)
(respectively, \(G'\)). Then \(T\) and \(T'\) are
isomorphic.}\footnote{Editor's note: This follows from \SGARefLink{sga3:xxiv:2.7}{2.7}, since
\(\underline{\operatorname{Isomext}}(G,G')\) has a section
(cf.\ the preceding editor's note).}

